% Source: https://artofproblemsolving.com/community/c3h2028579_circle_fact
% Stable identity: aops-topic-2028579-q01
(a) The point $\left(\dfrac{r}{\sqrt2},\dfrac{r}{\sqrt2}\right)$ lies on the circle, since $\dfrac{r^2}{2}+\dfrac{r^2}{2}=r^2$, but $x+y=r\sqrt2\ne r$. (Also $(-r,0)$ gives $x+y=-r$.) So the claim is false.

(b) If $x^2+y^2=r^2$ and $x+y=r$, then
\[r^2=(x+y)^2=x^2+y^2+2xy=r^2+2xy,\]
so $xy=0$. If $x=0$ then $y=r$; if $y=0$ then $x=r$. Both points work. So the line $x+y=r$ meets the circle in exactly two points,
\[\boxed{(r,0)\ \text{and}\ (0,r)}.\]
\begin{asy}
usepackage("amsmath");
size(5.5cm);
defaultpen(fontsize(10pt));
real r=2;
draw((-2.8,0)--(2.8,0),gray,Arrow(5)); draw((0,-2.8)--(0,2.8),gray,Arrow(5));
draw(circle((0,0),r),linewidth(0.9));
draw((-0.7,2.7)--(2.7,-0.7),blue);
dot((r,0),red); dot((0,r),red);
pair Q=r*dir(45); dot(Q);
label("$(r,0)$",(r,0),SW,red); label("$(0,r)$",(0,r),SW,red);
label("$\left(\dfrac{r}{\sqrt2},\dfrac{r}{\sqrt2}\right)$",Q,NE);
label("$x+y=r$",(2.7,-0.7),SE,blue);
label("$x^2+y^2=r^2$",r*dir(215),SW);
label("\mbox{}",(min(currentpicture,user=true)+max(currentpicture,user=true))/2,black); // keeps the text after the figure black when LaTeX typesets the labels inline
\end{asy}
